 >>?% ======================================================================
% CAPITULO 5 --- A Linguagem Geometrica: Gauges e Transformacoes
% Baseado em: Farinelli, S. (2021). arXiv:0910.1671v10, Secoes 2.2 e 2.3
%            Smith, A., Speed, C. (1998). Gauge Transforms in Stochastic Investment Modelling.
%            Bleecker, D. (1981). Gauge Theory and Variational Principles.
% ======================================================================
\chapter{A Linguagem Geom\'etrica: Gauges e Transforma\c{c}\~oes}

Este cap\'itulo constr\'oi a ponte entre finan\c{c}as e geometria
diferencial introduzindo o conceito central da Teoria Geom\'etrica da
Arbitragem (GAT): o \textbf{gauge} --- um par deflator--estrutura a
termo que codifica toda a informa\c{c}\~ao financeira de um instrumento.
Veremos que gauges formam um espa\c{c}o sobre o qual atua um grupo de
transforma\c{c}\~oes, exatamente como na f\'isica de part\'iculas o
potencial vetor $A_\mu$ se transforma sob $U(1)$. Esta analogia n\~ao
\'e metaf\'orica --- \'e matem\'atica. O fibrado principal do mercado
(Cap\'itulo~6) ser\'a constru\'ido inteiramente a partir dos objetos
definidos aqui.

A refer\^encia prim\'aria \'e Farinelli (2021), Se\c{c}\~oes~2.2 e~2.3
\cite{farinelli2021}. A formula\c{c}\~ao original de transforma\c{c}\~oes
de gauge em modelagem estoc\'astica de investimentos remonta a Smith e
Speed (1998) \cite{smith1998gauge}. A ponte com a teoria de gauge
diferencial segue Bleecker (1981) \cite{bleecker1981}.

% ======================================================================
\section{Por que Geometria?}
% ======================================================================

A f\'isica te\'orica do s\'eculo XX ensinou que as intera\c{c}\~oes
fundamentais da natureza s\~ao descritas por \textbf{teorias de gauge}:
o eletromagnetismo \'e uma teoria de gauge $U(1)$, a for\c{c}a fraca
\'e $SU(2)$, a for\c{c}a forte \'e $SU(3)$. Em cada caso, a
estrutura matem\'atica subjacente \'e um \textbf{fibrado principal}
$P \to M$ com grupo de estrutura $G$, munido de uma \textbf{conex\~ao}
$\omega$ cuja \textbf{curvatura} $F = d\omega + \omega \wedge \omega$
codifica o campo de for\c{c}a.

\begin{obs}[Analogia Eletromagnetismo $\leftrightarrow$ Mercado Financeiro]
O quadro abaixo estabelece a correspond\^encia precisa entre os dois
dom\'inios. N\~ao se trata de uma met\'afora: a estrutura alg\'ebrica e
geom\'etrica \'e rigorosamente a mesma.
\[
\begin{array}{c|c}
\textbf{Eletromagnetismo (F\'isica)} & \textbf{Mercado Financeiro (GAT)} \\ \hline
\text{Fibrado principal } P(M, U(1)) & \text{Fibrado principal } P(\mathcal{M}, G^*) \\
\text{Potencial vetor } A_\mu & \text{Conex\~ao de mercado } \omega \\
\text{Transforma\c{c}\~ao de gauge } A \mapsto A + d\Lambda & (D,P) \mapsto (D^\pi, P^\pi) \\
\text{Tensor de campo } F_{\mu\nu} = \partial_\mu A_\nu - \partial_\nu A_\mu & \text{Curvatura } F = d\omega + \omega \wedge \omega \\
\text{Carga el\'etrica } e & \text{Intensidade de cashflow } \pi \\
\text{Equa\c{c}\~oes de Maxwell } dF = 0,\; d*F = J & \text{Condi\c{c}\~ao de n\~ao-arbitragem } F = 0
\end{array}
\]
A consequ\^encia mais profunda: \textbf{arbitragem \'e curvatura}. Um
mercado livre de arbitragem \'e plano ($F = 0$); a presen\c{c}a de
oportunidades de arbitragem corresponde a curvatura n\~ao-nula na
conex\~ao de mercado.
\end{obs}

Por que esta linguagem \'e \'util em finan\c{c}as? Tr\^es raz\~oes:

\begin{enumerate}[(i)]
  \item \textbf{Invari\^ancia:} O pre\c{c}o de um derivativo n\~ao deve
        depender da escolha de numer\'ario. Exatamente como as leis da
        f\'isica n\~ao dependem da escolha de gauge, o valor de um
        contrato \'e \textbf{invariante de gauge} --- uma quantidade
        geom\'etrica global, n\~ao coordenada-dependente.
  \item \textbf{Unifica\c{c}\~ao:} Diferentes instrumentos financeiros
        (t\'itulos, a\c{c}\~oes, taxas de c\^ambio, derivativos) s\~ao
        todos descritos pelo mesmo objeto matem\'atico --- o gauge. As
        rela\c{c}\~oes entre eles s\~ao transforma\c{c}\~oes de gauge.
  \item \textbf{Classifica\c{c}\~ao:} A \^orbita de um gauge sob a
        a\c{c}\~ao do grupo $G^*$ determina sua \textbf{hierarquia de
        informa\c{c}\~ao}. Instrumentos em \^orbitas superiores
        cont\^em mais dados de mercado.
\end{enumerate}

% ======================================================================
\section{Deflatores e Estruturas a Termo}
% ======================================================================

\subsection{Defini\c{c}\~ao Formal}

\begin{defi}[Gauge]\label{def:gauge}
Seja $(\Omega, \mathcal{A}, \mathbb{A}, P)$ um espa\c{c}o de
probabilidade filtrado satisfazendo as condi\c{c}\~oes usuais
(Cap\'itulo~4). Um \textbf{gauge} \'e um par $(D, P)$ onde:
\begin{enumerate}[(i)]
  \item $D = (D_t)_{t \geq 0}$ \'e um processo estoc\'astico
        estritamente positivo, adaptado a $\mathbb{A}$, chamado
        \textbf{deflator}: o valor \textit{spot} (presente) do
        instrumento no instante $t$, expresso em unidades do
        numer\'ario escolhido.
  \item $P = (P_{t,s})_{0 \leq t \leq s}$ \'e um processo estoc\'astico
        a dois par\^ametros, estritamente positivo, $\mathcal{A}_t$-mensur\'avel
        para cada $s$ fixo, chamado \textbf{estrutura a termo}: o valor
        em $t$ de receber uma unidade do instrumento no instante
        $s \geq t$.
\end{enumerate}
com as propriedades fundamentais:
\[
P_{t,t} = 1, \qquad P_{t,s} > 0 \quad \text{para todo } 0 \leq t \leq s.
\]
\end{defi}

\begin{obs}[Interpreta\c{c}\~ao Econ\^omica]
Um gauge \'e a descri\c{c}\~ao completa e autocontida de um instrumento
financeiro. Ele responde a duas perguntas fundamentais:
\begin{itemize}[nosep]
  \item \textbf{Quanto vale hoje?} A resposta \'e $D_t$.
  \item \textbf{Quanto valer\'a no futuro, se eu travar o pre\c{c}o agora?}
        A resposta \'e $P_{t,s}$.
\end{itemize}
A condi\c{c}\~ao $P_{t,t} = 1$ \'e a normaliza\c{c}\~ao natural: uma
unidade entregue instantaneamente vale exatamente uma unidade. A
positividade $P_{t,s} > 0$ reflete o princ\'ipio econ\^omico de que
dinheiro futuro certo (em $s$) tem valor presente positivo (em $t$), por
menor que seja devido ao desconto.
\end{obs}

\subsection{Rela\c{c}\~ao com a Abordagem Cl\'assica}

Na abordagem cl\'assica (Cap\'itulo~4), o mercado \'e descrito por um
vetor de pre\c{c}os $S_t = (S_t^0, S_t^1, \dots, S_t^N)$ e uma conta
caixa $S_t^0 = \exp(\int_0^t r_u\,du)$. O gauge correspondente \'e:
\[
D_t = (S_t^0)^{-1} S_t^j, \qquad
P_{t,s} = \frac{\mathbb{E}_{\mathbb{Q}}[(S_s^0)^{-1} \mid \mathcal{A}_t]}
                 {(S_t^0)^{-1}}
\]
Ou seja, $D_t$ \'e o pre\c{c}o descontado do ativo $j$ e $P_{t,s}$
\'e o fator de desconto forward impl\'icito na medida neutra ao risco
$\mathbb{Q}$. A GAT, contudo, trabalha diretamente com $(D,P)$ sem
invocar $\mathbb{Q}$ --- a estrutura a termo \'e um objeto primitivo,
n\~ao derivado de uma medida martingal.

% ======================================================================
\section{Exemplos Concretos de Gauges}
% ======================================================================

Cada exemplo ilustra como instrumentos financeiros radicalmente
diferentes s\~ao unificados sob o mesmo formalismo. Em todos os casos,
o deflator $D_t$ e a estrutura a termo $P_{t,s}$ s\~ao expressos em
unidades do mesmo numer\'ario (tipicamente caixa, $S_t^0$).

\begin{ex}[T\'itulo de Cupom Zero (\textit{Zero-Coupon Bond})]\label{ex:zero_bond}
Uma fam\'ilia de t\'itulos p\'ublicos sem risco de cr\'edito que pagam
uma unidade monet\'aria no vencimento $T$.

\noindent\textbf{Deflator:} $D_t \equiv 1$. O valor do t\'itulo
emitido em $t$ \'e constante em unidades de caixa porque ele j\'a
representa dinheiro futuro descontado.

\noindent\textbf{Estrutura a termo:} $P_{t,s}$ \'e o pre\c{c}o em $t$
de um t\'itulo de cupom zero que vence em $s$. Tipicamente,
$P_{t,s} < 1$ para $s > t$ (o desconto reflete o valor temporal do
dinheiro). A fun\c{c}\~ao $s \mapsto P_{t,s}$ \'e chamada
\textbf{curva de desconto} no instante $t$.

\noindent\textbf{Propriedade:} Se o mercado \'e livre de arbitragem,
existe uma taxa curta $r_t$ tal que $P_{t,s} = \mathbb{E}_{\mathbb{Q}}[\exp(-\int_t^s
r_u\,du) \mid \mathcal{A}_t]$. Esta \'e a representa\c{c}\~ao
martingal cl\'assica --- a GAT n\~ao a exige, mas a permite.
\end{ex}

\begin{ex}[\'Indice de A\c{c}\~oes com Dividendos Reinvestidos]\label{ex:equity_index}
Considere um \'indice de a\c{c}\~oes (S\&P 500, Ibovespa) onde todos
os dividendos s\~ao automaticamente reinvestidos na compra de mais a\c{c}\~oes
do \'indice.

\noindent\textbf{Deflator:} $D_t =$ valor do \'indice em $t$,
descontado pela conta caixa. $D_t$ captura o \textbf{retorno total}
(pre\c{c}o + dividendos), descontado para eliminar o efeito da
infla\c{c}\~ao monet\'aria.

\noindent\textbf{Estrutura a termo:} $P_{t,s} =$ pre\c{c}o de um
contrato forward sobre o \'indice, emitido em $t$ com entrega em $s$,
expresso em unidades de $D_t$. Para a\c{c}\~oes, tipicamente $P_{t,s}$
\', e pr\'oximo de 1, pois a\c{c}\~oes n\~ao possuem `desconto
temporal'` no mesmo sentido que t\'itulos.

\noindent\textbf{Observa\c{c}\~ao:} O reinvestimento de dividendos \'e
crucial. Sem ele, $D_t$ refletiria apenas a aprecia\c{c}\~ao de
capital, e a estrutura a termo $P_{t,s}$ n\~ao seria bem definida
(pois haveria `vazamento'` de valor via dividendos em dinheiro).
\end{ex}

\begin{ex}[Taxas de C\^ambio como Quociente de Deflatores]\label{ex:fx}
Sejam $(D^{\text{dom}}, P^{\text{dom}})$ e $(D^{\text{for}},
P^{\text{for}})$ os gauges de dois pa\'ises (dom\'estico e estrangeiro).
A taxa de c\^ambio $X_t$ (unidades da moeda dom\'estica por unidade da
moeda estrangeira) \'e dada pelo quociente de deflatores:
\[
X_t = \frac{D_t^{\text{dom}}}{D_t^{\text{for}}}
\]
Esta rela\c{c}\~ao \'e uma consequ\^encia direta do princ\'ipio de
que deflatores convertem valores a numer\'ario. A estrutura a termo da
taxa de c\^ambio (forward FX) \'e obtida de forma an\'aloga:
\[
F_{t,s}^{\text{FX}} = X_t \cdot \frac{P_{t,s}^{\text{for}}}{P_{t,s}^{\text{dom}}}
\]
que \'e a conhecida \textbf{paridade coberta da taxa de juros}.
\end{ex}

\begin{ex}[Infla\c{c}\~ao]\label{ex:inflation}
O gauge de infla\c{c}\~ao descreve o poder de compra da moeda.

\noindent\textbf{Deflator:} $D_t = I_t$, o \'indice de pre\c{c}os ao
consumidor (IPC) no instante $t$. Ele mede quantas unidades monet\'arias
s\~ao necess\'arias para comprar uma cesta fixa de bens.

\noindent\textbf{Estrutura a termo:} $P_{t,s} =$ pre\c{c}o em $t$ de
um t\'itulo indexado a infla\c{c}\~ao que paga $I_s$ unidades
monet\'arias em $s$, expresso em unidades de $I_t$. Em outras
palavras, $P_{t,s}$ captura a taxa de juros real forward.
\end{ex}

% ======================================================================
\section{Taxa Forward Instant\^anea e Taxa Curta}
% ======================================================================

Suponha que a estrutura a termo $P_{t,s}$ seja diferenci\'avel em
$s$ para cada $t$ fixo. Isto \'e verdade para a maioria dos modelos
param\'etricos (Nelson--Siegel, Svensson) e para difus\~oes
markovianas.

\begin{defi}[Taxa Forward Instant\^anea]\label{def:forward_rate}
A \textbf{taxa forward instant\^anea} no instante $t$ para o
per\'iodo infinitesimal em $s \geq t$ \'e:
\[
\boxed{f_{t,s} := -\frac{\partial}{\partial s}\log P_{t,s}
       = -\frac{1}{P_{t,s}}\frac{\partial P_{t,s}}{\partial s}}
\]
\end{defi}

\begin{defi}[Taxa Curta (\textit{Short Rate})]\label{def:short_rate}
A \textbf{taxa curta} (ou taxa instant\^anea de juros) \'e o limite
da taxa forward quando o horizonte colapsa para o presente:
\[
\boxed{r_t := \lim_{s \to t^+} f_{t,s}
       = -\left.\frac{\partial}{\partial s}\log P_{t,s}\right|_{s=t}}
\]
\end{defi}

\begin{prop}[Reconstru\c{c}\~ao da Estrutura a Termo]\label{prop:reconstrucao}
A estrutura a termo pode ser reconstru\'ida a partir da taxa forward
via integra\c{c}\~ao:
\[
P_{t,s} = \exp\left(-\int_t^s f_{t,u}\,du\right)
\]
Em particular, se $r_u$ \'e a taxa curta e o mercado \'e livre de
arbitragem (implicando $f_{t,u} = \mathbb{E}_{\mathbb{Q}}[r_u \mid
\mathcal{A}_t]$ para um modelo de taxa curta), obtemos a f\'ormula
cl\'assica de desconto.
\end{prop}

\begin{proof}
Integrando a equa\c{c}\~ao diferencial $\frac{\partial}{\partial s}\log
P_{t,s} = -f_{t,s}$ de $s=t$ a $s=T$, com a condi\c{c}\~ao de contorno
$\log P_{t,t} = \log 1 = 0$:
\[
\log P_{t,T} - \log P_{t,t} = -\int_t^T f_{t,u}\,du
\implies
P_{t,T} = \exp\left(-\int_t^T f_{t,u}\,du\right)
\]
\end{proof}

\begin{obs}[Interpreta\c{c}\~ao Financeira]
A taxa forward $f_{t,s}$ \'e a taxa de juros marginal que o mercado
\"precifica\" hoje ($t$) para o per\'iodo $[s, s+ds]$. \'E an\'aloga
ao custo marginal em economia: o que custa emprestar dinheiro por
um instante adicional no futuro. A taxa curta $r_t$ \'e o custo
marginal do dinheiro \textit{agora}.
\end{obs}

% ======================================================================
\section{Condi\c{c}\~ao de Juros Positivos}
% ======================================================================

\begin{defi}[Condi\c{c}\~ao de \textit{Storage} (Juros Positivos)]\label{def:storage}
Um gauge $(D,P)$ satisfaz a \textbf{condi\c{c}\~ao de juros positivos}
(ou condi\c{c}\~ao de armazenamento) se, para todo $t$ e todo
$s_1 \leq s_2$:
\[
P_{t,s_1} \geq P_{t,s_2}
\]
isto \'e, a estrutura a termo \'e \textbf{mon\'otona decrescente} em
$s$. Equivalentemente, $f_{t,s} \geq 0$ para todo $s \geq t$.
\end{defi}

\begin{obs}[Significado Econ\^omico]
A condi\c{c}\~ao de juros positivos reflete o princ\'ipio do
\textbf{armazenamento} (\textit{storage}): dinheiro hoje \'e sempre
prefer\'ivel a dinheiro amanh\~a, pois pode ser armazenado sem custo
(em caixa) e mantido at\'e o futuro. Se $P_{t,s_1} < P_{t,s_2}$ para
$s_1 < s_2$, haveria uma oportunidade de arbitragem: comprar o
instrumento de prazo mais curto, receber o principal em $s_1$, guardar
o dinheiro f\'isico at\'e $s_2$, e lucrar a diferen\c{c}a.
\end{obs}

\begin{ex}[Instrumentos que \textbf{N\~AO} satisfazem a condi\c{c}\~ao de \textit{storage}]
\begin{enumerate}[(a)]
  \item \textbf{Infla\c{c}\~ao:} A estrutura a termo da infla\c{c}\~ao
        (Exemplo~\ref{ex:inflation}) tipicamente \'e \textbf{crescente}
        em $s$, pois $I_s$ tende a crescer com o tempo. N\~ao h\'a
        como ``armazenar'' o n\'ivel geral de pre\c{c}os --- a
        infla\c{c}\~ao \'e um processo n\~ao-armazen\'avel.
  \item \textbf{Dividendos:} O gauge de a\c{c}\~oes com dividendos
        (Exemplo~\ref{ex:equity_index}) tem $P_{t,s}$ pr\'oximo de~1,
        n\~ao monotonicamente decrescente. Empresas crescem e pagam
        dividendos; n\~ao \'e poss\'ivel ``armazenar'' o valor de uma
        a\c{c}\~ao sem risco.
  \item \textbf{Commodities com custo de carregamento:} Petr\'oleo,
        gr\~aos e metais t\^em custos de armazenamento f\'isico que
        podem tornar $P_{t,s}$ crescente em $s$ (contango).
\end{enumerate}
\end{ex}

A condi\c{c}\~ao de \textit{storage} n\~ao \'e uma exig\^encia do
formalismo --- \'e uma propriedade que alguns instrumentos possuem e
outros n\~ao. A GAT funciona igualmente para ambos, mas a estrutura
geom\'etrica \'e mais rica quando a condi\c{c}\~ao \'e satisfeita,
pois ent\~ao o semigrupo $G$ (Se\c{c}\~ao~\ref{sec:semigrupo}) tem
uma interpreta\c{c}\~ao probabil\'istica natural.

% ======================================================================
\section{Transforma\c{c}\~oes de Gauge}
% ======================================================================

\subsection{Motiva\c{c}\~ao}

Um investidor raramente det\'em um \'unico instrumento puro. Mais
comumente, ele possui um \textbf{portf\'olio de fluxos de caixa} ---
uma combina\c{c}\~ao de pagamentos em diferentes datas futuras. A GAT
formaliza esta ideia: dada uma ``receita'' de fluxos de caixa $\pi$,
constr\'oi-se um novo gauge a partir de um gauge base.

\begin{defi}[Transforma\c{c}\~ao de Gauge]\label{def:gauge_transform}
Seja $(D,P)$ um gauge e $\pi: [0,+\infty[ \to \mathbb{R}$ uma
fun\c{c}\~ao localmente integr\'avel chamada \textbf{intensidade de
fluxo de caixa} (\textit{cashflow intensity}). A transforma\c{c}\~ao de
gauge por $\pi$ produz um novo gauge $(D^\pi, P^\pi)$ definido por:
\[
\boxed{D_t^\pi := D_t \int_0^\infty \pi_h\, P_{t,t+h}\,dh}
\qquad
\boxed{P_{t,s}^\pi := \frac{\displaystyle\int_0^\infty \pi_h\, P_{t,s+h}\,dh}
                              {\displaystyle\int_0^\infty \pi_h\, P_{t,t+h}\,dh}}
\]
onde a interpreta\c{c}\~ao \'e:
\begin{itemize}[nosep]
  \item $\pi_h\,dh$ mede quanto de fluxo de caixa \'e recebido no
        horizonte $h$ (medido a partir de $t$);
  \item $P_{t,t+h}$ desconta esse fluxo de $t+h$ para $t$;
  \item $D_t^\pi$ \'e o valor presente total de todos os fluxos de
        caixa futuros, ponderados por $\pi$.
\end{itemize}
\end{defi}

\begin{obs}[Interpreta\c{c}\~ao Financeira]
Uma transforma\c{c}\~ao de gauge \textbf{constr\'oi um novo instrumento
a partir de um existente}. A fun\c{c}\~ao $\pi$ \'e a ``receita'' ---
ela especifica quanto de cada maturidade $h$ entra no novo
instrumento. O denominador em $P_{t,s}^\pi$ garante a
normaliza\c{c}\~ao $P_{t,t}^\pi = 1$.
\end{obs}

\subsection{Representa\c{c}\~ao com Medidas}

\'E conveniente interpretar $\pi$ como uma medida (assinada) na reta
real. Se $\pi$ cont\'em \textbf{fun\c{c}\~oes delta de Dirac}, podemos
representar pagamentos discretos. A nota\c{c}\~ao com integral de
Lebesgue--Stieltjes \'e:
\[
D_t^\pi := D_t \int_0^\infty P_{t,t+h}\,d\pi(h), \qquad
P_{t,s}^\pi := \frac{\int_0^\infty P_{t,s+h}\,d\pi(h)}
                       {\int_0^\infty P_{t,t+h}\,d\pi(h)}
\]
Esta formula\c{c}\~ao unifica pagamentos cont\'inuos e discretos.

\begin{ex}[T\'itulo com Cupom usando Deltas de Dirac]\label{ex:coupon_bond}
Para um t\'itulo com valor de face 1, taxa de cupom anual $g$, e
vencimento em $T$ anos (com pagamentos anuais):
\[
\pi = \sum_{k=1}^{T-1} g\,\delta_{k} + (1+g)\,\delta_{T}
\]
onde $\delta_a$ \'e a medida de Dirac concentrada em $a$. O primeiro
termo representa os $T-1$ cupons anuais de valor $g$; o segundo, o
principal mais o \'ultimo cupom no vencimento. Substituindo na
f\'ormula:
\begin{align*}
  D_t^\pi &= D_t\left[\sum_{k=1}^{T-1} g\,P_{t,t+k} + (1+g)P_{t,t+T}\right] \\[4pt]
  P_{t,s}^\pi &= \frac{\sum_{k=1}^{T-1} g\,P_{t,s+k} + (1+g)P_{t,s+T}}
                       {\sum_{k=1}^{T-1} g\,P_{t,t+k} + (1+g)P_{t,t+T}}
\end{align*}
\end{ex}

\begin{ex}[Perpetuidade (\textit{Consol})]\label{ex:consol}
Uma perpetuidade paga cupom cont\'inuo de 1 unidade por ano para
sempre. A fun\c{c}\~ao de cashflow \'e:
\[
\pi_h = 1 \quad \text{para todo } h \geq 0
\]
\'E a fun\c{c}\~ao de Heaviside ($\pi = \Theta$). O deflator da
perpetuidade \'e:
\[
D_t^{\text{consol}} = D_t \int_0^\infty P_{t,t+h}\,dh
\]
que \'e a \'area sob a curva de desconto.
\end{ex}

% ======================================================================
\section{O Semigrupo $G$ e o Grupo $G^*$}\label{sec:semigrupo}
% ======================================================================

A composi\c{c}\~ao de transforma\c{c}\~oes de gauge forma uma
estrutura alg\'ebrica not\'avel que ser\'a o grupo de estrutura do
fibrado principal do mercado (Cap\'itulo~6).

\begin{defi}[O Semigrupo $G$]\label{def:semi_grupo}
Seja $E'([0,+\infty[,\mathbb{R})$ o espa\c{c}o das distribui\c{c}\~oes
com suporte compacto em $[0,+\infty[$, munido da opera\c{c}\~ao de
\textbf{convolu\c{c}\~ao}:
\[
(\pi * \nu)(t) := \int_0^t \pi(t-u)\,d\nu(u)
\]
Ent\~ao $G := (E'([0,+\infty[,\mathbb{R}), *)$ \'e um \textbf{semigrupo
abeliano} com elemento neutro $\delta_0$ (a delta de Dirac na origem).
\end{defi}

\begin{defi}[O Grupo $G^*$]\label{def:grupo_estrela}
$G^* \subset G$ \'e o conjunto dos elementos invert\'iveis de $G$:
\[
G^* := \{\pi \in G \mid \exists\,\nu \in G \text{ tal que } \pi * \nu = \delta_0\}
\]
$G^*$ \'e um \textbf{grupo abeliano} sob convolu\c{c}\~ao.
\end{defi}

\begin{prop}[Propriedade Fundamental da Composi\c{c}\~ao]\label{prop:composicao}
A composi\c{c}\~ao de transforma\c{c}\~oes de gauge corresponde
\`a convolu\c{c}\~ao das intensidades de cashflow:
\[
\boxed{((D,P)^\pi)^\nu = (D,P)^{\pi * \nu}}
\]
Em particular, se $\pi \in G^*$ com inversa $\pi^{-1}$, ent\~ao
$(D,P)^\pi$ pode ser revertido ao gauge original via
$((D,P)^\pi)^{\pi^{-1}} = (D,P)$.
\end{prop}

\begin{proof}
Seja $(D^\pi, P^\pi)$ o gauge transformado. Aplicando a segunda
transforma\c{c}\~ao com intensidade $\nu$:
\begin{align*}
  D_t^{(\pi,\nu)}
  &= D_t^\pi \int_0^\infty \nu_h\, P_{t,t+h}^\pi\,dh \\[4pt]
  &= D_t \left(\int_0^\infty \pi_u\, P_{t,t+u}\,du\right) \cdot
     \frac{\int_0^\infty \nu_h \int_0^\infty \pi_u\, P_{t,t+h+u}\,du\,dh}
          {\int_0^\infty \pi_u\, P_{t,t+u}\,du} \\[4pt]
  &= D_t \int_0^\infty (\pi * \nu)_w\, P_{t,t+w}\,dw = D_t^{\pi * \nu}
\end{align*}
O c\'alculo para $P_{t,s}$ \'e an\'alogo, e a normaliza\c{c}\~ao
cancela o denominador.
\end{proof}

\begin{obs}[Estrutura de Semigrupo e N\~ao-negatividade]
Se nos restringirmos a medidas n\~ao-negativas, obtemos um
sub-semigrupo $G_+ \subset G$ cujos elementos t\^em interpreta\c{c}\~ao
financeira direta (todo fluxo de caixa recebido \'e positivo). O grupo
$G^*$ cont\'em elementos com partes negativas (ex.: transforma\c{c}\~ao
de taxa curta, Se\c{c}\~ao~\ref{sec:especiais}), correspondendo a
opera\c{c}\~oes financeiras que envolvem tanto recebimentos quanto
pagamentos.
\end{obs}

% ======================================================================
\section{\^Orbitas e Hierarquia de Informa\c{c}\~ao}
% ======================================================================

A a\c{c}\~ao do grupo $G^*$ sobre o espa\c{c}o de gauges particiona
este espa\c{c}o em \textbf{\^orbitas}. A estrutura das \^orbitas
codifica a hierarquia de informa\c{c}\~ao dispon\'ivel no mercado.

\begin{defi}[\^Orbita de um Gauge]\label{def:orbita}
A \textbf{\^orbita} de um gauge $(D,P)$ sob a a\c{c}\~ao de $G^*$ \'e:
\[
\mathcal{O}(D,P) := \{(D^\pi, P^\pi) \mid \pi \in G^*\}
\]
Dois gauges est\~ao na \textbf{mesma \^orbita} se e somente se existe
uma transforma\c{c}\~ao invert\'ivel mapeando um no outro.
\end{defi}

\begin{prop}[Hierarquia de Informa\c{c}\~ao]\label{prop:hierarquia}
Sejam $(D,P)$ e $(D^*,P^*)$ dois gauges.
\begin{enumerate}[(i)]
  \item Se $(D^*,P^*) = (D^\pi, P^\pi)$ com $\pi \in G^*$, ent\~ao
        ambos cont\^em a \textbf{mesma informa\c{c}\~ao de mercado}, e
        um pode ser reconstru\'ido a partir do outro.
  \item Se $(D^*,P^*)$ \'e obtido a partir de $(D,P)$ por uma
        transforma\c{c}\~ao \textbf{n\~ao-invert\'ivel}
        ($\pi \in G \setminus G^*$), ent\~ao a \^orbita de $(D^*,P^*)$
        \'e \textbf{estritamente inferior}: cont\'em menos
        informa\c{c}\~ao de mercado, e o gauge original n\~ao pode ser
        recuperado.
\end{enumerate}
\end{prop}

\begin{obs}[Intui\c{c}\~ao Geom\'etrica]
Imagine o espa\c{c}o de todos os gauges como um cone. As \^orbitas de
$G^*$ s\~ao ``fatias horizontais'' deste cone. Transforma\c{c}\~oes
invert\'iveis movem um gauge dentro de sua pr\'opria fatia (mesmo
n\'ivel de informa\c{c}\~ao). Transforma\c{c}\~oes singulares
projetam o gauge para uma fatia inferior (perda de informa\c{c}\~ao).
No topo do cone est\'a o \textbf{gauge principal} --- aquele a partir
do qual todos os outros podem ser obtidos por transforma\c{c}\~ao
(invert\'ivel ou n\~ao).
\end{obs}

\begin{ex}[Hierarquia na Pr\'atica]
Considere um mercado de t\'itulos:
\begin{itemize}[nosep]
  \item \textbf{Gauge de taxa curta} $(D^r, P^r)$: obtido do gauge
        de t\'itulos $(D^{\text{bond}}, P^{\text{bond}})$ pela
        transforma\c{c}\~ao $[-1]$ (Se\c{c}\~ao~\ref{sec:especiais}).
        Esta transforma\c{c}\~ao \'e \textbf{singular}, logo o gauge
        de taxa curta pertence a uma \^orbita inferior --- a partir
        dele n\~ao se pode reconstruir toda a estrutura a termo.
  \item \textbf{Gauge de t\'itulo com cupom}: obtido do gauge de
        zero-coupon por uma transforma\c{c}\~ao \textbf{invert\'ivel},
        logo pertence \`a mesma \^orbita. Ambos cont\^em a mesma
        informa\c{c}\~ao.
\end{itemize}
\end{ex}

% ======================================================================
\section{Transforma\c{c}\~oes Especiais}\label{sec:especiais}
% ======================================================================

Certas intensidades de cashflow $\pi$ t\^em significado financeiro
t\~ao fundamental que recebem nota\c{c}\~ao pr\'opria.

\begin{defi}[Transforma\c{c}\~oes Can\^onicas]\label{def:canonicas}
Definimos as seguintes intensidades can\^onicas:
\[
\begin{array}{ccl}
\text{Elemento neutro} & [0] &:= \delta_0 \\[4pt]
\text{Taxa curta (\textit{short rate})} & [-1] &:= \delta'_0\
  \text{(derivada da delta de Dirac)} \\[4pt]
\text{Perpetuidade (\textit{consol})} & [+1] &:= \Theta\
  \text{(fun\c{c}\~ao de Heaviside: $\Theta_h = 1$ para $h \geq 0$)}
\end{array}
\]
\end{defi}

\begin{prop}[\'Algebra dos \'Indices]\label{prop:algebra_indices}
Os \'indices can\^onicos satisfazem a seguinte \'algebra de
convolu\c{c}\~ao:
\[
\boxed{[k] * [l] = [k + l]}
\]
para todo $k, l \in \mathbb{Z}$.
\end{prop}

\begin{proof}
Por indu\c{c}\~ao a partir dos casos base:
\begin{itemize}[nosep]
  \item $[0] * [k] = [k]$ para todo $k$: pela defini\c{c}\~ao de
        elemento neutro.
  \item $[+1] * [+1] = [+2], [+1] * [+1] * [+1] = [+3]$, etc.: a
        convolu\c{c}\~ao de $n$ fun\c{c}\~oes de Heaviside produz
        $\Theta^{*n}(t) = \frac{t^{n-1}}{(n-1)!}$ para $t \geq 0$.
  \item $[-1] * [+1] = [0]$: a derivada da delta convolu\'ida com
        Heaviside devolve a delta ($\delta' * \Theta = \delta$).
  \item $[-1] * [-1] = [-2]$: segunda derivada da delta.
\end{itemize}
Estendendo para $k,l \in \mathbb{Z}$ arbitr\'arios, obtemos a
\'algebra de convolu\c{c}\~ao $[k]*[l] = [k+l]$. O grupo gerado por
$\{[k]\}_{k \in \mathbb{Z}}$ \'e isomorfo a $(\mathbb{Z}, +)$.
\end{proof}

\begin{obs}[Significado Financeiro dos \'Indices]
O \'indice $k$ conta quantas ``camadas de integra\c{c}\~ao'' ou
``camadas de deriva\c{c}\~ao'' separam o gauge do gauge de refer\^encia:
\begin{itemize}[nosep]
  \item $k = 0$: o pr\'oprio gauge (identidade).
  \item $k = +1$: perpetuidade --- integra a estrutura a termo sobre
        todos os horizontes (acumula\c{c}\~ao de fluxos).
  \item $k = -1$: taxa curta --- deriva a estrutura a termo na origem
        (extrai informa\c{c}\~ao marginal).
  \item $k = +n$ ($n > 1$): $n$-\'esima perpetuidade iterada.
  \item $k = -n$ ($n > 1$): $n$-\'esima derivada na origem.
\end{itemize}
\end{obs}

\begin{ex}[Construindo o Gauge de Taxa Curta]
Partindo do gauge de zero-coupon $(D^{\text{bond}}, P^{\text{bond}})$:
\[
D_t^{[-1]} = D_t \int_0^\infty \delta_0'(h)\, P_{t,t+h}\,dh
            = -D_t \left.\frac{\partial}{\partial h}P_{t,t+h}\right|_{h=0}
            = D_t \cdot r_t
\]
onde usamos $P_{t,t} = 1$ e $f_{t,t} = r_t$. A estrutura a termo
correspondente $P_{t,s}^{[-1]}$ \'e a raz\~ao de duas taxas forward,
mas o ponto crucial \'e que a transforma\c{c}\~ao $[-1]$ \'e
\textbf{singular}: $\delta_0'$ n\~ao possui inversa em $G$. Isto
reflete o fato econ\^omico de que conhecer apenas a taxa curta $r_t$
n\~ao basta para reconstruir toda a estrutura a termo --- \'e
necess\'ario conhecer a din\^amica completa de $r_t$ sob $\mathbb{Q}$.
\end{ex}

\begin{ex}[Gauge de Pre\c{c}o de A\c{c}\~ao a partir do Gauge de Dividendos]
Se $(D^{\text{div}}, P^{\text{div}})$ \'e o gauge de dividendos
(pagamentos de dividendos apenas, sem aprecia\c{c}\~ao de capital), o
gauge de pre\c{c}o da a\c{c}\~ao (com dividendos reinvestidos) \'e
obtido aplicando $[+1]$:
\[
D_t^{\text{acao}} = D_t^{\text{div}} \int_0^\infty P_{t,t+h}^{\text{div}}\,dh
\]
que \'e o valor presente de todos os dividendos futuros --- o modelo
de Gordon--Shapiro em linguagem de gauge.
\end{ex}

% ======================================================================
\section{Diagrama de Gauges}
% ======================================================================

A figura a seguir (\`a la diagrama comutativo) resume as
rela\c{c}\~oes entre os principais tipos de gauges e as
transforma\c{c}\~oes que os conectam. As setas s\~ao rotuladas pela
intensidade de cashflow $\pi$ (ou pelo \'indice $[k]$) que realiza a
transforma\c{c}\~ao.

\[
\begin{array}{ccccccc}
& & & \boxed{\text{Principal } (D,P)} & & & \\[4pt]
& & \swarrow[-1] & \downarrow[+1] & \searrow\text{Cupom} & & \\[4pt]
& \boxed{\text{Taxa Curta}} & & \boxed{\text{Consol}} & & \boxed{\text{Coupon Bond}} & \\[4pt]
& \swarrow[-1] & & \downarrow[+1] & & \swarrow[-1] & \\[4pt]
\boxed{\text{Volatilidade}} & & \boxed{\text{Renda Fixa}} & & \boxed{\text{Fwd Rate}} & & \boxed{\text{Zero Bond}} \\[4pt]
& & & \downarrow[+1] & & & \\[4pt]
& & & \boxed{\text{Equity}} & & & \\[4pt]
& & & \downarrow\text{Div} & & & \\[4pt]
& & & \boxed{\text{Dividend}} & & &
\end{array}
\]

\begin{center}
\small
\textbf{Legenda:} $[+1]$ = perpetuidade (Heaviside $\Theta$);
$[-1]$ = taxa curta ($\delta'$). As setas pontilhadas indicam
transforma\c{c}\~oes singulares (perda de informa\c{c}\~ao: mudam de
\^orbita). As setas s\'olidas indicam transforma\c{c}\~oes invert\'iveis
(mesma \^orbita).
\end{center}

\begin{obs}[Leitura do Diagrama]
\begin{itemize}[nosep]
  \item O \textbf{gauge principal} (topo) cont\'em a m\'axima
        informa\c{c}\~ao. A partir dele, qualquer outro gauge pode
        ser obtido.
  \item Descendo no diagrama via setas pontilhadas ($[-1]$), perde-se
        informa\c{c}\~ao (proje\c{c}\~ao para \^orbita inferior).
  \item Movimentos horizontais (ex.: Principal $\to$ Coupon Bond)
        s\~ao transforma\c{c}\~oes invert\'iveis --- mesma \^orbita,
        mesma informa\c{c}\~ao.
  \item O gauge de \textbf{equity} (a\c{c}\~ao com dividendos
        reinvestidos) \'e obtido aplicando $[+1]$ ao gauge de renda
        fixa, enquanto o gauge de \textbf{dividend} (apenas dividendos)
        \'e obtido removendo a aprecia\c{c}\~ao de capital.
\end{itemize}
\end{obs}

% ======================================================================
\section{S\'intese e Conex\~ao com os Pr\'oximos Cap\'itulos}
% ======================================================================

O cap\'itulo estabeleceu o vocabul\'ario fundamental da GAT:

\begin{enumerate}[(i)]
  \item \textbf{Gauge} $(D,P)$: a ``certid\~ao de nascimento''
        financeira de qualquer instrumento.
  \item \textbf{Transforma\c{c}\~ao de gauge} via $\pi$: constru\c{c}\~ao
        de novos instrumentos a partir de existentes.
  \item \textbf{Grupo $G^*$}: as transforma\c{c}\~oes invert\'iveis,
        que particionam o espa\c{c}o de gauges em \^orbitas.
  \item \textbf{\^Orbitas e hierarquia}: cada \^orbita corresponde a
        um n\'ivel de informa\c{c}\~ao de mercado.
  \item \textbf{\'Algebra $[k]*[l] = [k+l]$}: aritm\'etica das
        transforma\c{c}\~oes fundamentais.
\end{enumerate}

No Cap\'itulo~6, veremos que o conjunto de todos os gauges forma um
\textbf{fibrado principal} $P(\mathcal{M}, G^*)$ sobre o espa\c{c}o de
par\^ametros de mercado $\mathcal{M}$. O grupo $G^*$ atua como grupo
de estrutura, e as transforma\c{c}\~oes de gauge s\~ao exatamente as
mudan\c{c}as de trivializa\c{c}\~ao local. Esta \'e a ponte
definitiva entre finan\c{c}as e geometria diferencial.

\vspace{12pt}
\begin{center}
\rule{0.4\textwidth}{0.5pt}\\[4pt]
{\small \textbf{Princ\'ipio Geom\'etrico:} O que a f\'isica chama de
\emph{simetria de gauge}, as finan\c{c}as chamam de \emph{mudan\c{c}a
de numer\'ario}. A invari\^ancia da f\'isica sob $U(1)$ \'e a
invari\^ancia das finan\c{c}as sob $G^*$.}\\[4pt]
\rule{0.4\textwidth}{0.5pt}
\end{center}


